\section{Conclusion}
Digital payment use was common in the recruited sample. The migrant accounts show how paying through a relative's wallet can solve an immediate problem without giving the user control of the account. A transfer within Colombia can also succeed even when someone else must arrange delivery abroad. Who can act after a failure depends on which step went wrong.

The survey distinguishes recent domestic use from what respondents report doing with incoming remittances. The observed patterns also distinguish reported onboarding conditions from the discretion and information women say they have. The exploratory profiles illustrate why reported transfer ability, balance holding, and discretion need not move together; they do not establish user types. A count of completed payments does not tell us whether women could pay again without help, whether recipients could use the money, or whether users could obtain a remedy.

A useful evaluation would follow payment attempts through delivery or failure, including attempts that are abandoned. It would ask whether customers keep access when a number or document changes, compare the cost and timing of the options they can actually use, and distinguish a complaint, a response and recovery. Transaction counts alone cannot answer those questions. We do not estimate the effect of Bre-B or a change in welfare.
